\documentclass{article}
\title{MC911 - Projeto 1}

\begin{document}

\includegraphics{unicamp}

\maketitle

\textbf{\textit{Introduction}}

For the first project, it was proposed the creation and implementation of a grammar, which is supposed to take as input a LaTex file, usually used to create papers and scientific articles, and then, convert them to HTML format as output, s used to create and display of web pages.

To do this, it was created a LALR(1) grammar, taking a finite automata that uses a stack to parse. The implementation was done using Flex and Bison tools (an lexical analyser and parser respectively). Not all of the LaTeX commands had effect in the HTML scenario, so, some of then should be just ignored and others interpreted.

\textbf{\textit{Usage}}

To use, you need to compile first. Just run:

	make

After that, just execute the parser binary. One note, all the html code comes from the stdout, while erros, warnings and info messagens come from stderr. So it is best to execute with the following scheme:

	./parser < manual.tex > manual.html 

\textit{\textbf{Supported Commands}}

Follows the list of implemented commands with all their expected results in HTML:

\begin{itemize}

	\item \textbf{Header Commands}
		\begin{itemize}
			 \item documentclass: recognized, but not effective
			 \item usepackage:  recognized, but not effective
			 \item title: define the title of the LaTex document, so it could be displayed with \textbf{maketitle} command;
			 \item author: recognized, but not effective
		\end{itemize}

	\item \textbf{Body Commands:}
		\begin{itemize}
	 		\item \textit{begin}: marks the begining of the document;
	 		\item \textit{end:} marks the ending of the document;
	 		\item \textit{maketitle: mostra o} shows the title defined in \"title\" command in the header;
	 		\item \textit{textbf:} apply bold to texto1;
	 		\item \textit{textit}: apply italic to texto1;
	 		\item \textit{begin}: generate an enumerated list of itens (treat nested list);
	 		\item \textit{item texto1:} includes texto1 as an item inside the itens list; 
	 		\item \textit{end}: marks the end of the item list;
	 		\item \textit{includegraphics:} shows figura1;
	 		\item \textit{cite:} add the number of the ref that is in the section \'thebibliography\';
	 		\item \textit{begin:} begin the section of bibliography as a list of references;
	 		\item \textit{bibitem:} define a reference item;
	 		\item \textit{end:} add the references to the body of the document (one per document);
	 		\item \textit{\$:} the text between two dollars serves to mark the math environment;
		\end{itemize}

\end{itemize}

\textit{\textbf{The Math Text}}
To deal with the math mode, we used MathJax library. Using an two-phase processing of the HTML text, the library check the content looking for makers \$ that define the math environment, and then, it substitute of the text in LaTex format to the correct Math format and formulas.

For example, the following text is a math text:

"When $a \ne 0$, there are two solutions to $ax^2 + bx + c = 0$ and they are $x = {-b \pm \sqrt{b^2-4ac} \over 2a}.$"

\textit{\textbf{The Grammar}}

In LaTex, despite some commands follow an order, in general the text that must be printed in the final document is mixed with the commands that format this same text (for instance, commands like bold os italic). We attempted to do a grammar that is simple enough to do not have conflicts, but ,at same time, powerfull enough to analyse commands, text and math mode (using \$).

To guide the process, we used some principles of construction grammar described in \cite{Appel02} and \cite{Aho86}.

latex $\rightarrow$ preamble document

preamble $\rightarrow$ DOCCLASS LBRACE normalt RBRACE headerlist 
preamble $\rightarrow$ DOCCLASS LBRACKET normalt RBRACKET LBRACE normalt RBRACE headerlist

headerlist $\rightarrow$ headerlist USEPKG LBRACE headertext RBRACE
headerlist $\rightarrow$ headerlist USEPKG LBRACKET headertext RBRACKET LBRACE headertext RBRACE
headerlist $\rightarrow$ TITLE LBRACE headertext RBRACE
headerlist $\rightarrow$ AUTHOR LBRACE headertext RBRACE
headerlist $\rightarrow$ headerlist NEWLINES
headerlist $\rightarrow$

document $\rightarrow$ BEGINDOC body ENDDOC
document $\rightarrow$ BEGINDOC ENDDOC

body $\rightarrow$ body text
body $\rightarrow$ body command
body $\rightarrow$ text
body $\rightarrow$ command

command $\rightarrow$ INGRAPH LBRACE normalt RBRACE
command $\rightarrow$ startbib biblist ENDBIB

startbib $\rightarrow$ BEGINBIB

biblist $\rightarrow$  BBITEM LBRACE normalt RBRACE headertext
biblist $\rightarrow$ BBITEM LBRACE normalt RBRACE headertext
biblist $\rightarrow$ NEWLINES
biblist $\rightarrow$ biblist NEWLINES

specialsymbol $\rightarrow$ BEGINITEM
specialsymbol $\rightarrow$ ENDITEM
specialsymbol $\rightarrow$ ITEM
specialsymbol $\rightarrow$ MAKETITLE
specialsymbol $\rightarrow$ CITE LBRACE normalt RBRACE

headertext $\rightarrow$ normalt
headertext $\rightarrow$ italict
headertext $\rightarrow$ boldt
headertext $\rightarrow$ headertext normalt
headertext $\rightarrow$ headertext italict
headertext $\rightarrow$ headertext boldt

text $\rightarrow$ normalt
text $\rightarrow$ italict
text $\rightarrow$ boldt
text $\rightarrow$ NEWLINES
text $\rightarrow$ math

math $\rightarrow$ DOLLAR mathexp DOLLAR

mathexp $\rightarrow$ mathexp STRING
mathexp $\rightarrow$ mathexp CHAR
mathexp $\rightarrow$ mathexp LBRACE
mathexp $\rightarrow$ mathexp RBRACE
mathexp $\rightarrow$ STRING
mathexp $\rightarrow$ CHAR
mathexp $\rightarrow$ LBRACE
mathexp $\rightarrow$ RBRACE

normalt $\rightarrow$ STRING
normalt $\rightarrow$ CHAR
normalt $\rightarrow$ specialsymbol
normalt $\rightarrow$ LBRACE
normalt $\rightarrow$ RBRACE

boldt $\rightarrow$ TXTBF LBRACE boldtext RBRACE

italict $\rightarrow$ TXTIT LBRACE italictext RBRACE

boldtext $\rightarrow$ CHAR
boldtext $\rightarrow$ STRING
boldtext $\rightarrow$ specialsymbol
boldtext $\rightarrow$ italict
boldtext $\rightarrow$ NEWLINES
boldtext $\rightarrow$ boldtext CHAR
boldtext $\rightarrow$ boldtext STRING
boldtext $\rightarrow$ boldtext specialsymbol
boldtext $\rightarrow$ boldtext italict
boldtext $\rightarrow$ boldtext NEWLINES

italictext $\rightarrow$ CHAR
italictext $\rightarrow$ STRING
italictext $\rightarrow$ specialsymbol
italictext $\rightarrow$ italict
italictext $\rightarrow$ NEWLINES
italictext $\rightarrow$ italictext CHAR
italictext $\rightarrow$ italictext STRING
italictext $\rightarrow$ italictext specialsymbol
italictext $\rightarrow$ italictext italict
italictext $\rightarrow$ italictext NEWLINES


\begin{thebibliography}

\bibitem{Appel02} Appel, A., \textbf{Modern Compiler Implementation in Java, 2nd Ed.}, October 21, 2002

\bibitem{Aho86} Aho, A., Sethi, M., Ullman, J., \textbf{Compilers: Principles, Techniques and Tools}, 1986

\end{thebibliography}



\end{document}
